%%%%%%%%%%%%%%%%%
% This is an sample CV template created using altacv.cls
% (v1.7.2, 28 August 2024) written by LianTze Lim (liantze@gmail.com). Compiles with pdfLaTeX, XeLaTeX and LuaLaTeX.
%
%% It may be distributed and/or modified under the
%% conditions of the LaTeX Project Public License, either version 1.3
%% of this license or (at your option) any later version.
%% The latest version of this license is in
%%    http://www.latex-project.org/lppl.txt
%% and version 1.3 or later is part of all distributions of LaTeX
%% version 2003/12/01 or later.
%%%%%%%%%%%%%%%%

%% Use the "normalphoto" option if you want a normal photo instead of cropped to a circle
% \documentclass[10pt,a4paper,normalphoto]{altacv}

\documentclass[8pt,a4paper,ragged2e,withhyper]{altacv}
%% AltaCV uses the fontawesome5 and simpleicons packages.
%% See http://texdoc.net/pkg/fontawesome5 and http://texdoc.net/pkg/simpleicons for full list of symbols.
%EXPERIENCE
%Senior researcher, CNES : LISA esa/NASA Space mission
% Change the page layout if you need to
\geometry{left=1.25cm,right=1.25cm,top=.5cm,bottom=1.cm,columnsep=1.2cm}

% The paracol package lets you typeset columns of text in parallel
\usepackage{paracol}
\usepackage{fontawesome5}
\usepackage{hyperref}

% Change the font if you want to, depending on whether
% you're using pdflatex or xelatex/lualatex
% WHEN COMPILING WITH XELATEX PLEASE USE
% xelatex -shell-escape -output-driver="xdvipdfmx -z 0" sample.tex
\iftutex 
  % If using xelatex or lualatex:
  \setmainfont{Roboto Slab}
  \setsansfont{Lato}
  \renewcommand{\familydefault}{\sfdefault}
\else
  % If using pdflatex:
  \usepackage[rm]{roboto}
  \usepackage[defaultsans]{lato}
  % \usepackage{sourcesanspro}
  \renewcommand{\familydefault}{\sfdefault}
\fi

% Change the colours if you want to
\definecolor{SlateGrey}{HTML}{2E2E2E}
\definecolor{LightGrey}{HTML}{666666}
\definecolor{DarkPastelRed}{HTML}{8AC0D9}%{450808}
\definecolor{PastelRed}{HTML}{00B0DA}%{8F0D0D}
\definecolor{GoldenEarth}{HTML}{B4AD0C}%{E7D192}
\colorlet{name}{black}
\colorlet{tagline}{PastelRed}
\colorlet{heading}{DarkPastelRed}
\colorlet{headingrule}{GoldenEarth}
\colorlet{subheading}{PastelRed}
\colorlet{accent}{PastelRed}
\colorlet{emphasis}{SlateGrey}
\colorlet{body}{LightGrey}

\definecolor{violet}{HTML}{5A2582}
\definecolor{rouge}{HTML}{F53B3B}
\definecolor{noir}{HTML}{00232C}
\definecolor{bleu}{HTML}{00B0DA}
\definecolor{rose}{HTML}{F831A0}
\definecolor{jaune}{HTML}{FFEC00}
\definecolor{vert}{HTML}{34AD6C}


% Change some fonts, if necessary
\renewcommand{\namefont}{\Huge\rmfamily\bfseries}
\renewcommand{\personalinfofont}{\footnotesize}
\renewcommand{\cvsectionfont}{\LARGE\rmfamily\bfseries}
\renewcommand{\cvsubsectionfont}{\large\bfseries}
%\newcommand{\orcid}[1]{\href{https://orcid.org/#1}{\includesvg[width=10pt]{orcid}}}

% Change the bullets for itemize and rating marker
% for \cvskill if you want to
\renewcommand{\cvItemMarker}{{\small\textbullet}}
\renewcommand{\cvRatingMarker}{\faCircle}
% ...and the markers for the date/location for \cvevent
% \renewcommand{\cvDateMarker}{\faCalendar*[regular]}
% \renewcommand{\cvLocationMarker}{\faMapMarker*}


% If your CV/résumé is in a language other than English,
% then you probably want to change these so that when you
% copy-paste from the PDF or run pdftotext, the location
% and date marker icons for \cvevent will paste as correct
% translations. For example Spanish:
% \renewcommand{\locationname}{Ubicación}
% \renewcommand{\datename}{Fecha}


%% Use (and optionally edit if necessary) this .tex if you
%% want to use an author-year reference style like APA(6)
%% for your publication list
% \input{pubs-authoryear.tex}

%% Use (and optionally edit if necessary) this .tex if you
%% want an originally numerical reference style like IEEE
%% for your publication list
\input{pubs-num.tex}

%% sample.bib contains your publications
\addbibresource{sample.bib}
% \usepackage{academicons}\let\faOrcid\aiOrcid
\begin{document}
\name{BOILEAU Guillaume} %, PhD
\tagline{Systems Test \& Validation Engineer | Integration, Technical Analysis \& Complex Software Systems}%Your Position or Tagline Here}
%% You can add multiple photos on the left or right
\photoL{2.8cm}{moi}
%\photoR{2.5cm}{Suitcase_High}%6892(1)}

\personalinfo{%
  % Not all of these are required!
  \email{guillaume.boileaupro@gmail.com}
  \phone{+33 6 59 94 85 84}
  \location{Alpes Maritimes, FRANCE}
  %\mailaddress{1 avenue des citronniers 06800 Cagnes sur Mer, France}
  %\homepage{https://www.researchgate.net/profile/Guillaume-Boileau-2}%\href{https://www.researchgate.net/profile/Guillaume-Boileau-2}{Guillaume-Boileau}}
 %\faLinkedin \href{https://www.linkedin.com/in/guillaume-boileau-phd-891b81265/}{in/guillaume-boileau}
     

  % \twitter{@twitterhandle}
  %\xtwitter{@x-handle}
  \linkedin{guillaume-boileau-phd-891b81265}
  %\github{your_id}
  \researchgate{Guillaume-Boileau-2}
  \orcid{0000-0002-3576-6968}
  
  %% You can add your own arbitrary detail with
  %% \printinfo{symbol}{detail}[optional hyperlink prefix]
  % \printinfo{\faPaw}{Hey ho!}[https://example.com/]

  %% Or you can declare your own field with
  %% \NewInfoFiled{fieldname}{symbol}[optional hyperlink prefix] and use it:
  % \NewInfoField{gitlab}{\faGitlab}[https://gitlab.com/]
  % \gitlab{your_id}
  %%
  %% For services and platforms like Mastodon where there isn't a
  %% straightforward relation between the user ID/nickname and the hyperlink,
  %% you can use \printinfo directly e.g.
  % \printinfo{\faMastodon}{@username@instace}[https://instance.url/@username]
  %% But if you absolutely want to create new dedicated info fields for
  %% such platforms, then use \NewInfoField* with a star:
  % \NewInfoField*{mastodon}{\faMastodon}
  %% then you can use \mastodon, with TWO arguments where the 2nd argument is
  %% the full hyperlink.
  % \mastodon{@username@instance}{https://instance.url/@username}
}

\makecvheader
%% Depending on your tastes, you may want to make fonts of itemize environments slightly smaller

\vspace{-0.3cm}
\AtBeginEnvironment{itemize}{\small}

%% Set the left/right column width ratio to 6:4.
\columnratio{0.64}

% Start a 2-column paracol. Both the left and right columns will automatically
% break across pages if things get too long.
\begin{paracol}{2}
\cvsection{Experience}

\cvevent{\textbf{Software Engineer – System Integration \& Hardware Interfaces}}{VEV Platform Services France}{2025 -- 2026}{Valbonne, France}

\textbf{Context:} Development and integration of software services for EV charging infrastructure within a CPMS platform.

\textbf{Objective:} Deliver reliable software components, operational data flows and robust service interactions with charging station systems.

\begin{itemize}
\item \textbf{Software development:} Developed web and backend services using TypeScript, Angular and Node.js.
\item \textbf{System integration:} Contributed to the integration and maintenance of software services interacting with EV charging infrastructure.
\item \textbf{Operational validation:} Supported service verification and monitoring to ensure consistency and reliability of operational data exchanges.
\item \textbf{Technical issue analysis:} Contributed to the investigation and resolution of software and hardware interaction issues in production environments.
\item \textbf{Data flows:} Implemented and monitored FTP-based services for operational data exchange.
\end{itemize}

\divider

\cvevent{\textbf{Senior R\&D Engineer – Complex Software Systems, Integration \& Validation}}{Laboratoire Lagrange, Observatoire de la Côte d'Azur, CNES}{2023 -- 2025}{Nice, France}

\textbf{Context:} Engineering activities for the LISA space mission, involving complex software chains, analysis pipelines and coordinated technical developments in an international framework.

\textbf{Objective:} Develop, structure and validate software and analysis workflows for complex signal-processing systems.

\begin{itemize}
\item \textbf{Software platform development:} Developed a Python-based platform for large-scale model integration, comparison and analysis workflows. \textcolor{DarkPastelRed}{\href{https://gitlab.oca.eu/gboileau/lisa_l3_tools}{\bf GitLab:CATpyLISA}}
\item \textbf{Integration and validation support:} Contributed to the consistency analysis and verification of complex software components and processing chains.
\item \textbf{Technical coordination:} Coordinated software development tasks, aligned contributions and organized collaborative technical workflows.
\item \textbf{Test and analysis campaigns:} Structured analysis scenarios, comparison campaigns and result validation workflows in a constrained collaborative environment.
\item \textbf{Technical issue investigation:} Investigated inconsistencies, analysed unexpected behaviours and supported technical issue resolution with contributors from different domains.
\item \textbf{Team management:} Supervised interns and supported engineering activities across technical tasks.
\end{itemize}

\divider

\cvevent{\textbf{Data Scientist – Noise Modeling for Precision Instruments}}{Universiteit Antwerpen, Belgium}{2021 -- 2023}{Antwerp, Belgium}

\textbf{Context:} Data analysis and performance studies for precision instruments in a high-requirement technical environment.

\textbf{Objective:} Identify, characterize and mitigate sources affecting system sensitivity using robust statistical and computational methods.

\begin{itemize}
\item \textbf{Analysis and validation workflows:} Developed robust analysis pipelines for noise source identification, characterization and mitigation.
\item \textbf{Performance investigation:} Contributed to validation activities and technical investigations related to system sensitivity limitations.
\item \textbf{Environmental analysis:} Studied environmental effects and their impact on measurement quality and system-level performance.
\item \textbf{Technical rigor:} Worked in an environment requiring traceability, reliability of results and structured technical workflows.
\end{itemize}

\divider

\cvevent{\textbf{Junior R\&D Engineer – Signal Analysis, Simulation \& Technical Investigations}}{ARTEMIS, Observatoire de la Côte d'Azur, Nice, France}{2018--2021}{Nice, France}

\textbf{Context:} Engineering and analysis activities for gravitational-wave mission preparation, involving signal analysis, simulation and noise characterization.

\textbf{Objective:} Develop robust analysis methods for complex signals and support technical investigations in demanding system environments.

\begin{itemize}
\item \textbf{Signal analysis:} Developed methods for the separation of overlapping low-SNR signals in complex datasets.
\item \textbf{Simulation studies:} Built simulation approaches for complex signal environments and large-scale datasets.
\item \textbf{Technical investigations:} Investigated noise sources through cross-sensor correlation analyses in diagnostic contexts.
\item \textbf{Technical contribution:} Contributed to collaborative developments related to analysis methods and software tools for complex space-system applications.
\end{itemize}




\switchcolumn

\iffalse
\cvsection{My Life Philosophy}

\begin{quote}
``Life is like riding a bicycle. To keep your balance, you must keep moving.''
\end{quote}


\cvsection{Summary and Strengths}

Experienced astrophysicist with advanced skills in statistical analysis, Bayesian modeling, and machine learning. Proven ability to lead teams in international collaborations. Strong scientific background demonstrated by impactful publications and presentations at major conferences. Skilled in managing large-scale scientific projects and solving complex problems.
\fi

\cvsection{Summary and Strengths}

\textbf{Engineer experienced in complex software and signal-processing environments}, with a strong background in \textbf{system integration}, \textbf{validation}, \textbf{technical analysis} and \textbf{robust software workflows}. Used to working on demanding systems requiring \textbf{rigor}, \textbf{traceability}, \textbf{structured execution} and \textbf{cross-functional collaboration}.

Experience in \textbf{technical investigations}, \textbf{result consistency analysis}, \textbf{automation-oriented workflows}, and coordination of contributors in complex engineering environments.

Additional hands-on experience in \textbf{software development and system interfaces} for EV charging infrastructure using TypeScript, Angular and Node.js, including operational data exchange services and production issue analysis.

\cvsection{Team Management}

\begin{itemize}
  \item Supervision of \textbf{3 Master’s thesis interns (M2)} during engineering and development phases.
  \item Coordination of technical activities involving \textbf{research engineers and technicians} on a \textbf{data processing pipeline} for the \textbf{LISA space mission}.
  \item Experience in organizing tasks, supporting contributors and maintaining alignment across collaborative technical developments in demanding environments.
\end{itemize}

\cvsection{Languages}
%\cvskill{French}{5}
%\divider
\cvskill{English}{4}
\divider

\cvskill{Spanish}{2}
\divider

%\cvskill{German}{3.5} %% Supports X.5 values.

%% Yeah I didn't spend too much time making all the
%% spacing consistent... sorry. Use \smallskip, \medskip,
%% \bigskip, \vspace etc to make adjustments.
\medskip

\cvsection{Education}
\cvevent{PhD in Physics}{Université Côte d'Azur}{2018 -- 2021}{Nice, France}
\divider


\cvevent{Selective Graduate Program in Fundamental Physics (Magistère)}{Université Paris-Saclay}{2016 -- 2018}{Orsay, France}
%\textit{with distinction}

\divider
\cvevent{MSc in Astronomy and Astrophysics}{Observatoire de Paris / Université Paris-Saclay}{2017 -- 2018}{Paris, France}
%\textit{with distinction}
\divider
\cvevent{BSc in Fundamental Physics}{Université Paris-Saclay}{2015 -- 2016}{Orsay, France}
\cvsection{Professional Training}

\cvevent{Supervised Machine Learning (Regression \& Classification)}{DeepLearning.AI – Andrew Ng}{2020}{Coursera}
%\textit{with distinction}

%\divider

%\cvevent{CPGE Classe préparatoire aux grandes écoles MPSI/MP (Maths-Physics)}{Lycée  Dessaigne}{2013-2015}{Blois, France}
%\textit{with distinction and merit}

%\divider

%\cvevent{High School Diploma, Baccalauréat  S }{Lycée  Augustin  Thierry}{2013}{Blois, France}

%\cventry{2013-2015}{CPGE Classe préparatoire aux grandes écoles MPSI/MP (Maths-Physics)}{Lycée  Dessaigne, Blois, France}{%(a  two  years preparatory  course  that  is  a  prerequisite  for  entry  into  one  of  the Grandes Ecoles)
%}{}{\textit{with distinction and merit}}

%\cventry{2013}{High School Diploma, Baccalauréat  S %(high  school diploma focusing on Mathematics and Physics) 
%}{Lycée  Augustin  Thierry}{Blois, France}{}{\textit{with distinction }}


\end{paracol}

\newpage
\cvsection{Projects}

\cvevent{Co-author and full member of the LISA Consortium}{ESA/NASA Space Mission -- Large-scale International Collaboration}{2018--now}{}

Contribution to software and analysis workflows for the LISA mission within an international engineering and scientific collaboration involving ESA, NASA and multiple research institutes.

\textbf{Role:} Contribution to software workflows, technical coordination, complex signal-processing activities and collaborative engineering tasks.

\textbf{Program scale:} Major international space mission with an estimated budget of approximately 2 billion EUR over 10 years.

\divider

\cvevent{Co-author and full member of Einstein Telescope and LIGO-Virgo-KAGRA Collaborations}{International Collaborations}{2021--now}{}

Contribution to collaborative developments in data analysis, signal characterization and software methods within the LIGO--Virgo--KAGRA and Einstein Telescope collaborations.

\textbf{Role:} Development of analysis tools, contribution to technical activities and support to complex software workflows in high-standard collaborative environments.

\textbf{Environment:} Large-scale international collaborations involving strong technical requirements, coordination and structured methods.

%\medskip

\cvsection{Strengths}
\vspace{-.3cm}
\cvsubsection{Core skills}
{\small
\cvtag{System Integration}
\cvtag{Verification \& Validation}
\cvtag{System Verification}
\cvtag{System Validation}
\cvtag{Integration Testing}
\cvtag{Test Execution}
\cvtag{Test Procedures}
\cvtag{Test Automation}
\cvtag{Technical Investigation}
\cvtag{Anomaly Investigation}
\cvtag{Root Cause Analysis}
\cvtag{Technical Diagnostics}
\cvtag{Issue Resolution}
\cvtag{Validation Strategy Support}
\cvtag{Test Data Management}
\cvtag{System-Level Analysis}
\cvtag{Complex Systems}
\cvtag{Signal Processing}
\cvtag{Software Integration}
\cvtag{Traceability}
\cvtag{Structured Execution}
\cvtag{Technical Coordination}
}
\\
\divider\smallskip
\vspace{-.3cm}

\cvsubsection{Technical skills}
{\small
\cvtag{Python} \cvtag{Git} \cvtag{Docker} \cvtag{Linux}
\cvtag{SQL} \cvtag{TypeScript} \cvtag{Angular} \cvtag{Node.js}
\cvtag{C} \cvtag{Matlab} \cvtag{Pandas}
\cvtag{REST API} \cvtag{bash scripting} \cvtag{FTP}
\cvtag{Issue Tracking} 

\cvtag{Technical Reporting}
\cvtag{Condor} \cvtag{Slurm} \cvtag{PyTorch}
\cvtag{MongoDB} \cvtag{Mathematica}\cvtag{Test Automation Scripting} \cvtag{CI/CD} \cvtag{Issue Tracking}
}
\\
\divider\smallskip
\vspace{-.3cm}

{\small\LaTeXraggedright
\cvsubsection{Soft skills}
{\small
\cvtag{Autonomy}
\cvtag{Analytical Thinking}
\cvtag{Execution Rigor}
\cvtag{Problem Solving}
\cvtag{Adaptability}
\cvtag{Responsibility}
\cvtag{Attention to Detail}
\cvtag{Collaboration}
\cvtag{Stakeholder Communication}
\cvtag{Structured Communication}
\cvtag{Team Coordination}
\cvtag{Customer Awareness}
\par}}

\cvtag{\textit{Applications: complex software systems, signal analysis, system integration, validation workflows}}
\end{document}
